\documentclass[11pt]{article}
\usepackage[a4paper,margin=1in]{geometry}
\usepackage[T1]{fontenc}
\usepackage[utf8]{inputenc}
\usepackage{lmodern}
\usepackage{amsmath,amssymb,amsthm,mathtools}
\usepackage{enumitem}
\usepackage{xcolor}
\usepackage[hidelinks]{hyperref}
\usepackage{microtype}

\newtheorem{theorem}{Theorem}[section]
\newtheorem{proposition}[theorem]{Proposition}
\newtheorem{lemma}[theorem]{Lemma}
\newtheorem{corollary}[theorem]{Corollary}
\theoremstyle{remark}
\newtheorem{remark}[theorem]{Remark}
\newtheorem{definition}[theorem]{Definition}

\newcommand{\kk}{\mathbf C}
\newcommand{\PP}{\mathbf P}
\newcommand{\m}{\mathfrak m}
\newcommand{\cO}{\mathcal O}
\newcommand{\length}{\operatorname{length}}
\newcommand{\Osc}{\operatorname{Osc}}
\newcommand{\gr}{\operatorname{gr}}

\hypersetup{
  pdftitle={Prescribed Tjurina Algebras and Complete Spectra in Osculating-Absorbing Gauss Fibres},
  pdfauthor={Lluis Eriksson}
}

\title{Prescribed Tjurina Algebras and Complete Spectra\\
in Osculating-Absorbing Gauss Fibres}
\author{Lluis Eriksson}
\date{August 25, 2026}

\begin{document}
\maketitle

\begin{abstract}
We prove a simultaneous realization theorem for point-span
osculating-absorbing Gauss fibres in every dimension.  Given the smallest
permitted support and any list of isolated tangent-section germs of order
at least \(s+1\), a smooth hypersurface of every degree
\[
 D\ge1+\sum_i(\tau_i+2)
\]
can be chosen whose fibre has exactly the prescribed completed Tjurina
algebras.  Its length and dual multiplicity are respectively
\(\sum_i\tau_i\) and \(\sum_i\mu_i\).  For surfaces, the classical complete spectrum
of ordinary plane multiple points then realizes every integer fibre length
between
\[
 \binom{m+2}{2}\left\lfloor\frac{3s^2+4s-3}{4}\right\rfloor
 \quad\text{and}\quad
 \binom{m+2}{2}s^2,
\]
while the dual multiplicity stays fixed.

We also give a short strong-Lefschetz construction of extremal
three-variable Tjurina algebras.  The explicit ordinary germs
\[
 x^{s+1}+y^{s+1}+z^{s+1}+(x+y+z)^{s+2}
\]
attain the value \(s(s+2)(2s-1)/3\), classically computed by Wahl as the
minimum for positive-weight deformations of the Fermat cone.  Our proof
combines an Euler-reduced universal lower bound with an initial-ideal
comparison and Stanley's theorem for
\(\mathbf C[x,y,z]/(x^s,y^s,z^s)\).  This yields sharp absorbed
threefold fibres at minimal reduced support.  All statements are over
\(\mathbf C\); the degree thresholds are sufficient and not claimed
minimal.
\end{abstract}

\section{Introduction}

For an isolated hypersurface germ
\(h\in\mathbf C[[x_1,\ldots,x_d]]\), its Tjurina and Milnor numbers are
\[
 \tau(h)=\dim_{\mathbf C}\frac{\mathbf C[[x_1,\ldots,x_d]]}
 {(h,\partial_1h,\ldots,\partial_dh)},
 \qquad
 \mu(h)=\dim_{\mathbf C}\frac{\mathbf C[[x_1,\ldots,x_d]]}
 {(\partial_1h,\ldots,\partial_dh)}.
\]
Fixing the order of \(h\) leaves a nontrivial minimization problem for
\(\tau\).  A recent finite-jet argument gives a universal family of lower
bounds and determines the exact plane minimum
\cite{ErikssonEuler2026}.  In three variables its specialization is
\[
 \tau(h)\ge \frac{s(s+2)(2s-1)}3,
\]
which is sharp.  For \(s\ge3\), the same value was already computed by
Wahl as the minimum on positive-weight deformations of the degree-\(s+1\)
Fermat cone \cite[Example~4.7]{Wahl1985}; see also
\cite[Example~4.6]{Almiron2022}.  We supply a short proof for the single
uniform extremizer obtained by adding the \((s+2)\)-nd power of a
Lefschetz linear form.  This is an explicit realization and a bridge to the
global application, not a priority claim for the numerical minimum.

The main purpose is geometric.  For a smooth hypersurface
\(X^d\subset\PP^{d+1}\), the completed local algebra of a Gauss fibre at a
point is the Tjurina algebra of the corresponding tangent-section germ
\cite{ErikssonGauss2026}.  Point-span order-\(s\) osculating absorption in
the complete \(m\)-uple embedding forces both order at least \(s+1\) and at
least \(\binom{d+m}{d}\) reduced support points
\cite{ErikssonOsculating2026,ErikssonEuler2026}.  Combining the universal
floor with the explicit Lefschetz realization gives a sharp
dimension-three fibre endpoint.

For surfaces, the endpoint was already known.  We strengthen it in a
different direction.  Classical work on ordinary plane singularities says
that every integral Tjurina number between the sharp minimum and the Milnor
number occurs \cite{BGM1988,CaninoEtAl2025,CaninoEtAlArxiv}.  We prove that
arbitrary finite lists of these germs can be installed simultaneously on an
extremal absorbed Gauss fibre.  Hence the entire integer interval of fibre
lengths occurs while the multiplicity of the dual hypersurface is fixed.

No exhaustive priority assertion is intended.  The local plane spectrum,
Wahl's three-variable minimum, and the strong Lefschetz input are
established results.  The content here is the prescribed-list realization
under the absorption constraint, its complete global spectrum consequence,
and the direct Lefschetz extremizer used for the threefold endpoint.

\section{The Euler-reduced lower bound in three variables}

We use the convention \(\binom ab=0\) for \(a<b\), including negative
upper indices.

\begin{lemma}[Euler reduction]\label{lem:euler}
Let \(R=\mathbf C[[x_1,\ldots,x_d]]\) and
\(h\in\m^{s+1}\).  Put
\[
 r=h-\frac1{s+1}\sum_{i=1}^d x_i\partial_i h.
\]
Then \(r\in\m^{s+2}\) and
\[
 (h,\partial_1h,\ldots,\partial_dh)
 =(r,\partial_1h,\ldots,\partial_dh).
\]
\end{lemma}

\begin{proof}
Euler's identity cancels the homogeneous component of degree \(s+1\).
The equality of ideals follows directly from the displayed definition of
\(r\).  Notice that the derivative generators are not replaced by the
derivatives of \(r\).
\end{proof}

\begin{theorem}[Exact three-variable floor]\label{thm:local3}
Let \(s\ge1\), and let \(h\in\mathbf C[[x,y,z]]\) have an isolated critical
point and order at least \(s+1\).  Then
\begin{equation}\label{eq:threefloor}
 \tau(h)\ge T^{(3)}_s:=\frac{s(s+2)(2s-1)}3.
\end{equation}
This is the best possible bound for every \(s\).
\end{theorem}

\begin{remark}[Historical scope]\label{rem:wahl}
For \(s\ge3\), Wahl computed the same sharp number in his study of
positive-weight deformations of
\(x^{s+1}+y^{s+1}+z^{s+1}\)
\cite[Example~4.7]{Wahl1985}; Almir\'on records the formula explicitly in
\cite[Example~4.6]{Almiron2022}.  The lower-bound proof below applies to
every isolated germ of order at least \(s+1\).  Proposition~\ref{prop:equality}
then gives a particularly uniform one-term extremizer.  We do not claim
priority for the numerical minimum or for the abstract existence of
extremal positive-weight deformations.
\end{remark}

\begin{proof}[Proof of the lower bound]
Work in \(R/\m^{2s}\), whose dimension is \(\binom{2s+2}{3}\).
By Lemma~\ref{lem:euler}, the Tjurina ideal has three generators of order at
least \(s\) and one generator of order at least \(s+2\).  Their images are
therefore spanned by at most
\[
 3\binom{s+2}{3}+\binom{s}{3}
\]
coefficient multiples.  Rank is at most the dimension of the source, so
\[
 \tau(h)\ge
 \binom{2s+2}{3}-3\binom{s+2}{3}-\binom{s}{3}
 =\frac{s(s+2)(2s-1)}3.
\]
The Tjurina algebra surjects onto this truncated quotient.  This proves the
claim, including \(s=1,2\) under the binomial convention.
\end{proof}

\section{A strong-Lefschetz extremizer}

Put \(\ell=x+y+z\) and
\begin{equation}\label{eq:extremizer}
 f_s=x^{s+1}+y^{s+1}+z^{s+1}+\ell^{s+2}.
\end{equation}
Its initial form is Fermat and defines a smooth plane curve.  Thus \(f_s\)
is an ordinary isolated hypersurface singularity of multiplicity \(s+1\),
and \(\mu(f_s)=s^3\).

\begin{lemma}[Lefschetz quotient]\label{lem:lefschetz}
For \(s\ge1\), let
\[
 A_s=\mathbf C[x,y,z]/(x^s,y^s,z^s).
\]
Then
\[
 \dim_{\mathbf C} A_s/(\ell^{s+2})
 =\frac{s(s+2)(2s-1)}3.
\]
\end{lemma}

\begin{proof}
For \(s=1\), the assertion is immediate from \(A_1=\mathbf C\), so assume
\(s\ge2\).
The linear form \(\ell=x+y+z\) is a strong Lefschetz element of the
monomial complete intersection \(A_s\)
\cite{Stanley1980,PhuongTran2023}.  If \(H_j=\dim_{\mathbf C}(A_s)_j\),
multiplication by \(\ell^{s+2}\) from degree \(j-s-2\) to degree \(j\)
has maximal rank.  Therefore
\[
 \dim\bigl(A_s/(\ell^{s+2})\bigr)_j
 =\max\{H_j-H_{j-s-2},0\}.
\]
The Hilbert function of \(A_s\) is symmetric and unimodal with socle degree
\(3(s-1)\).  Its positive difference
\(H_j-H_{j-s-2}\) occurs precisely for \(0\le j\le2s-1\), and the
difference is nonpositive thereafter (eventually both terms vanish).
Equivalently, the midpoint between \(j\) and \(j-s-2\) crosses the symmetry
centre between \(2s-1\) and \(2s\).  Summing the positive part through degree
\(2s-1\), using
\(\sum H_jt^j=(1-t^s)^3/(1-t)^3\), gives
\[
 \sum_{j=0}^{2s-1}(H_j-H_{j-s-2})
 =\binom{2s+2}{3}-3\binom{s+2}{3}-\binom{s}{3}
 =\frac{s(s+2)(2s-1)}3.
\]
\end{proof}

\begin{proposition}[Equality family]\label{prop:equality}
For every \(s\ge1\), the germ \(f_s\) in \eqref{eq:extremizer} satisfies
\[
 \mu(f_s)=s^3,
 \qquad
 \tau(f_s)=\frac{s(s+2)(2s-1)}3.
\]
\end{proposition}

\begin{proof}
Euler reduction applied to \eqref{eq:extremizer} gives
\[
 f_s-\frac1{s+1}(x\partial_xf_s+y\partial_yf_s+z\partial_zf_s)
 =-\frac1{s+1}\ell^{s+2}.
\]
Hence \(\ell^{s+2}\) belongs to the Tjurina ideal \(I_s\).  The lowest
homogeneous forms of its three derivative generators are nonzero multiples
of \(x^s,y^s,z^s\).  Thus the \(\m\)-adic initial ideal satisfies
\[
 \operatorname{in}_{\m}(I_s)
 \supseteq (x^s,y^s,z^s,\ell^{s+2}).
\]
Here \(\operatorname{in}_{\m}(I_s)\) denotes the ideal generated by the
lowest nonzero homogeneous forms of all elements of \(I_s\).  Length is
preserved on passage to the associated graded algebra, so
Lemma~\ref{lem:lefschetz} yields
\[
 \tau(f_s)\le\dim_{\mathbf C}A_s/(\ell^{s+2})=T^{(3)}_s.
\]
The reverse inequality is Theorem~\ref{thm:local3}.  Finally, the ordinary
leading form gives \(\mu(f_s)=s^3\).
\end{proof}

The containment of initial ideals, rather than an asserted equality, is
important.  It supplies only the upper bound needed to meet the universal
lower bound.

\section{Absorbed Gauss fibres}

Let \(X^d\subset\PP^{d+1}_{\mathbf C}\) be a smooth non-linear hypersurface,
let \(\gamma_X:X\to X^\vee\) be its Gauss morphism, and fix a tangent
hyperplane \(\eta=[W]\).  Write \(\Gamma_\eta\) for the scheme-theoretic
fibre and \(Z=(\Gamma_\eta)_{\mathrm{red}}\).  Let \(H=\cO_X(1)\) and
\[
 S_Z=\operatorname{Im}\left(
 H^0(Z,H^m|_Z)^*\longrightarrow H^0(X,H^m)^*\right).
\]
The support is \emph{point-span order-\(s\) osculating-absorbing} if
\[
 \widehat\Osc^s_p(H^m)\subseteq S_Z
 \qquad(p\in Z).
\]

We recall two inputs, with their scopes.  First, if \(h_p\) is a local
equation of \(X\cap W\) at \(p\), then
\begin{equation}\label{eq:fibrealgebra}
 \widehat{\cO}_{\Gamma_\eta,p}
 \cong\mathbf C[[z_1,\ldots,z_d]]/(h_p,\partial h_p).
\end{equation}
Second, absorption in the complete \(m\)-uple embedding, with
\(1\le s\le m\), implies
\begin{equation}\label{eq:support}
 h_p\in\m_p^{s+1},
 \qquad
 |Z|\ge N_{d,m}:=\binom{d+m}{d}.
\end{equation}
These facts are proved in
\cite{ErikssonGauss2026,ErikssonOsculating2026,ErikssonEuler2026}; they are
not reproved as new results here.

\begin{corollary}[Sharp numerical floor in dimension three]\label{cor:gauss3}
Suppose \(d=3\) and \(Z\) is point-span order-\(s\)
osculating-absorbing in the complete \(H^m\) embedding, where
\(1\le s\le m\).  Then
\[
 \length(\Gamma_\eta)
 \ge \frac{s(s+2)(2s-1)}3\,|Z|
 \ge \frac{s(s+2)(2s-1)}3\binom{m+3}{3}.
\]
Both factors can be attained simultaneously at sufficiently large ambient
degree, as proved in Theorem~\ref{thm:global3} below.
\end{corollary}

\begin{proof}
Sum Theorem~\ref{thm:local3} over \eqref{eq:fibrealgebra}, then apply
\eqref{eq:support}.
\end{proof}

\section{Simultaneous realization}

The following theorem isolates the interpolation argument used for both
geometric consequences.

General deformation-theoretic realization results for prescribed
singularities have a substantial literature, including
\cite{GreuelKarras1989}.  The theorem below uses an elementary
high-degree separator construction because it must retain the common
tangent hyperplane, unisolvent support, and absorption constraint.
Its separator--Bertini architecture already appears for two ordinary local
types in \cite{ErikssonGauss2026}; here we record the heterogeneous
finite-determinacy version and do not claim a new general method for
prescribed singularities.

\begin{theorem}[Prescribed Tjurina algebras on an extremal fibre]
\label{thm:prescribe}
Let \(d\ge1\), \(1\le s\le m\), and
\(N=\binom{d+m}{d}\).  Choose \(N\) isolated convergent hypersurface germs
\[
 g_i\in\mathbf C\{z_1,\ldots,z_d\},
 \qquad \operatorname{ord}(g_i)\ge s+1,
\]
and put
\[
 E=\sum_{i=1}^N(\tau(g_i)+2).
\]
If
\begin{equation}\label{eq:degree}
 D\ge E+1,
\end{equation}
then there exist a smooth integral hypersurface
\(X\subset\PP^{d+1}\) of degree \(D\), a hyperplane \(W\), and a tangent
hyperplane point \(\eta=[W]\in X^\vee\) such that:
\begin{enumerate}[label=(\roman*)]
\item \(Z=(\gamma_X^{-1}(\eta))_{\mathrm{red}}\) consists of exactly \(N\)
points and \(\dim S_Z=N<h^0(X,H^m)\);
\item \(Z\) is point-span order-\(s\) osculating-absorbing;
\item the tangent-section germ at the \(i\)-th point is contact equivalent
to \(g_i\), hence
\[
 \widehat{\cO}_{\Gamma_\eta,p_i}
 \cong\mathbf C[[z_1,\ldots,z_d]]/(g_i,\partial g_i);
\]
\item
\[
 \length(\Gamma_\eta)=\sum_{i=1}^N\tau(g_i),
 \qquad
 \operatorname{mult}_{\eta}(X^\vee)=\sum_{i=1}^N\mu(g_i).
\]
\end{enumerate}
The degree bound \eqref{eq:degree} is sufficient and is not asserted to be
minimal.
\end{theorem}

\begin{proof}
Choose \(N\) points \(Z\subset W\cong\PP^d\) for which degree-\(m\)
evaluation is an isomorphism.  One explicit choice in an affine chart is
the simplex lattice
\[
 [1:i_1:\cdots:i_d],
 \qquad i_j\in\mathbf Z_{\ge0},\quad \sum_j i_j\le m;
\]
multivariate Newton interpolation proves unisolvence.

Put \(a_i=\tau(g_i)+1\).  Isolated hypersurface germs are contact
\((\tau+1)\)-determined \cite[Corollary~2.24]{GLS2007}; prescribing the
\(a_i\)-jet therefore fixes the contact class of \(g_i\).  The elementary
separator product makes
\[
 H^0(\PP^d,\cO(D))\longrightarrow
 \bigoplus_{i=1}^N\cO_{\PP^d,p_i}/\m_{p_i}^{a_i+1}
\]
surjective once
\[
 D\ge\sum_i(a_i+1)-1=E-1.
\]
Indeed, an \(a_i\)-jet at \(p_i\) can be multiplied by
\(\prod_{j\ne i}\ell_{ij}^{a_j+1}\), where \(\ell_{ij}\) vanishes at
\(p_j\) and is a unit at \(p_i\).  The stronger bound
\eqref{eq:degree} leaves one degree after the product
\[
 P_u=\prod_{i=1}^N\ell_{i,u}^{a_i+1},
 \qquad \deg P_u=E,
\]
so the affine system with the prescribed jets generates first jets at
every \(u\in W\setminus Z\).

A member singular at a specified \(u\in W\setminus Z\) imposes \(d+1\)
linear conditions, while \(\dim W=d\).  Hence a general degree-\(D\)
equation \(f\) on \(W\) is smooth away from \(Z\) and has exactly the chosen
contact types at \(Z\).

Embed \(W=V(y)\) in \(\PP^{d+1}\) and set
\[
 F_G=f+yG,
 \qquad G\in H^0(\PP^{d+1},\cO(D-1)).
\]
Choose \(G\) general and nonzero at \(Z\).  Then \(X=V(F_G)\) is smooth: at
\(Z\) its normal derivative is \(G(p)\ne0\); on \(W\setminus Z\) this
follows from \(df\ne0\); and off \(W\) the singular-point incidence has
fibre codimension \(d+2>\dim\PP^{d+1}=d+1\).  The smooth hypersurface is
connected and therefore integral.  Locally on \(X\), \(y=-f/G\), so the
tangent-section germs are contact equivalent to the prescribed germs.

The singular locus of \(X\cap W\) is exactly \(Z\), hence the reduced Gauss
fibre over \([W]\) is exactly \(Z\).  Each \(\tau(g_i)\ge1\), while
\(N=\binom{d+m}{d}\ge m+1\); hence \(D\ge E+1>m\), and degree-\(m\) ambient
forms identify with \(H^0(X,H^m)\).  A form vanishing on \(Z\) restricts to
zero on \(W\) by unisolvence, and is consequently divisible by \(y\).
Along \(X\), \(y=-f/G\) has order at least \(s+1\) at every support point.  The
reverse-annihilator criterion gives order-\(s\) absorption, while evaluation
gives \(\dim S_Z=N\).  Properness follows from
\(h^0(X,H^m)=\binom{d+m+1}{d+1}>N\).

Finally, \eqref{eq:fibrealgebra} gives the length formula.  The dual
multiplicity formula is the sum of the local Milnor numbers; see
\cite{Parusinski1991} for the classical dual-multiplicity result and
\cite{ErikssonGauss2026,ErikssonDual2026} for its use in this fibre
setting.
\end{proof}

\section{Complete surface spectra and sharp threefold fibres}

Set
\[
 T^{(2)}_s=\left\lfloor\frac{3s^2+4s-3}{4}\right\rfloor.
\]
An ordinary plane germ of multiplicity \(s+1\) has Milnor number \(s^2\),
and its Tjurina number assumes every integer from \(T^{(2)}_s\) through
\(s^2\) \cite[Theorem~5.1]{CaninoEtAlArxiv}; that theorem in turn attributes
the spectrum to the classical sources recorded there.  We use this local
result rather than claiming it.

\begin{theorem}[Complete absorbed surface spectrum]\label{thm:spectrum2}
Let \(1\le s\le m\), put \(N=\binom{m+2}{2}\), and assume
\[
 D\ge(s^2+2)N+1.
\]
For every integer \(L\) satisfying
\[
 NT^{(2)}_s\le L\le Ns^2,
\]
there exist a smooth integral surface \(X\subset\PP^3\) of degree \(D\)
and a point-span order-\(s\) osculating-absorbing Gauss fibre with
\[
 |Z|=N,
 \qquad \length(\Gamma_\eta)=L,
 \qquad \operatorname{mult}_{\eta}(X^\vee)=Ns^2.
\]
More strongly, Theorem~\ref{thm:prescribe} realizes any specified list
of \(N\) ordinary plane germs of multiplicity \(s+1\), including their
completed Tjurina algebras.
\end{theorem}

\begin{proof}
Write \(L-NT^{(2)}_s=r_1+\cdots+r_N\) with
\(0\le r_i\le s^2-T^{(2)}_s\).  Such a decomposition exists by a greedy
choice.  For each \(i\), choose an ordinary plane germ with Tjurina number
\(T^{(2)}_s+r_i\), using the cited complete local spectrum, and apply
Theorem~\ref{thm:prescribe} with \(d=2\).
\end{proof}

\begin{theorem}[Sharp absorbed threefold fibres]\label{thm:global3}
Let \(1\le s\le m\), put
\[
 N=\binom{m+3}{3},
 \qquad T^{(3)}_s=\frac{s(s+2)(2s-1)}3,
\]
and assume
\[
 D\ge\bigl(T^{(3)}_s+2\bigr)N+1.
\]
Then there exist a smooth integral threefold
\(X\subset\PP^4\) of degree \(D\) and a point-span order-\(s\)
osculating-absorbing Gauss fibre such that
\[
 |Z|=N,
 \qquad
 \length(\Gamma_\eta)=T^{(3)}_sN,
 \qquad
 \operatorname{mult}_{\eta}(X^\vee)=s^3N.
\]
Thus Corollary~\ref{cor:gauss3} is sharp simultaneously in local length and
reduced support size.
\end{theorem}

\begin{proof}
Apply Theorem~\ref{thm:prescribe} with \(d=3\) and all prescribed
germs equal to \(f_s\).  Proposition~\ref{prop:equality} supplies their
Tjurina and Milnor numbers.
\end{proof}

\section{Scope, reproducibility, and open directions}

The argument is over \(\mathbf C\).  Euler division by \(s+1\), contact
determinacy, the stated Gauss-fibre inputs, and strong Lefschetz are used in
that setting.  No positive-characteristic extension is asserted.

The replay accompanying the paper has three deliberately separated layers.
Exact binomial arithmetic checks the formula and maximizing truncation for
\(1\le s\le50\).  Exact rational matrix ranks check the relevant Lefschetz
maps for \(1\le s\le8\).  Sparse quotient computations over two large prime
fields directly check the displayed germ for \(2\le s\le7\).  These finite
fixtures test the implementation and examples; they do not prove the
universal theorems.

The paper identifies the Euler-reduced floor with Wahl's sharp value in
three variables and realizes it by the displayed Lefschetz family; no
analogous assertion is made in four or more variables.  Classification of
all equality germs is also open.
For higher dimensions a recent, non-refereed preprint gives stronger
asymptotic ratio estimates and Fermat-deformation families
\cite{MaZuo2026}; it is not used here.
The global degree thresholds come from transparent separator products and
are probably far from optimal.  The surface spectrum fixes only numerical
length and dual multiplicity unless the stronger prescribed-list statement
is invoked; different Tjurina algebras of the same length need not be
isomorphic.

\section{Conclusion}

The prescribed-list theorem turns local Tjurina data into exact absorbed
Gauss fibres at minimal reduced support.  In the plane case it upgrades
endpoint sharpness to the full integer spectrum at fixed dual multiplicity.
In three variables, Euler reduction and strong Lefschetz give a direct
uniform representative of Wahl's sharp value and hence an exact global
endpoint for threefold hypersurfaces.  Equality classification,
higher-variable sharpness, and optimal ambient degrees remain distinct
problems.

\begin{thebibliography}{99}

\bibitem{Almiron2022}
P.~Almir\'on,
\emph{On the quotient of Milnor and Tjurina numbers for two-dimensional
isolated hypersurface singularities},
Math. Nachr. \textbf{295} (2022), 1254--1263.
\href{https://doi.org/10.1002/mana.202100371}
{doi:10.1002/mana.202100371}.

\bibitem{BGM1988}
J.~Brian\c{c}on, M.~Granger, and P.~Maisonobe,
\emph{Le nombre de modules du germe de courbe plane \(x^a+y^b=0\)},
Math. Ann. \textbf{279} (1988), 535--551.
\href{https://doi.org/10.1007/BF01456286}{doi:10.1007/BF01456286}.

\bibitem{CaninoEtAl2025}
S.~Canino, A.~Gimigliano, and M.~Id\`a,
\emph{On the Jacobian scheme of a plane curve},
Comm. Algebra \textbf{53} (2025), 582--592.
\href{https://doi.org/10.1080/00927872.2024.2384056}
{doi:10.1080/00927872.2024.2384056}.

\bibitem{CaninoEtAlArxiv}
S.~Canino, A.~Gimigliano, and M.~Id\`a,
\emph{On the Jacobian scheme of a plane curve}, arXiv:2302.07042v2 (2024),
especially Theorem~5.1 and Remark~5.8.
\url{https://arxiv.org/abs/2302.07042}.

\bibitem{ErikssonGauss2026}
L.~Eriksson,
\emph{Fat Gauss Fibres and Tjurina--Milnor Defects Forced by Osculating
Absorption}, ARR-2026-2MZWECWVEN97ARVQ, v1 (2026).
\url{https://arr-research.github.io/papers/ARR-2026-2MZWECWVEN97ARVQ/}.

\bibitem{ErikssonDual2026}
L.~Eriksson,
\emph{Exact multiplicity floors for dual singularities from absorbing Gauss
fibres}, ARR-2026-0WAPCGQHNC82S8VJ, v1 (2026).
\url{https://arr-research.github.io/papers/ARR-2026-0WAPCGQHNC82S8VJ/}.

\bibitem{ErikssonOsculating2026}
L.~Eriksson,
\emph{Exact Higher-Osculating Rank Floors over Arbitrary Fields},
ARR-2026-66Q8M61AA196T8BC, v2 (2026).
\url{https://arr-research.github.io/papers/ARR-2026-66Q8M61AA196T8BC/versions/v2/}.

\bibitem{ErikssonEuler2026}
L.~Eriksson,
\emph{Euler-Reduced Tjurina Floors for Osculating-Absorbing Gauss Fibres},
ARR-2026-5XEX8EX0629R997Y, v2 (2026).
\url{https://arr-research.github.io/papers/ARR-2026-5XEX8EX0629R997Y/versions/v2/}.

\bibitem{GLS2007}
G.-M.~Greuel, C.~Lossen, and E.~Shustin,
\emph{Introduction to Singularities and Deformations},
Springer Monographs in Mathematics, Springer, 2007.
\href{https://doi.org/10.1007/3-540-28419-2}{doi:10.1007/3-540-28419-2}.

\bibitem{GreuelKarras1989}
G.-M.~Greuel and U.~Karras,
\emph{Families of varieties with prescribed singularities},
Compos. Math. \textbf{69} (1989), 83--110.
\url{https://www.numdam.org/item/CM_1989__69_1_83_0/}.

\bibitem{MaZuo2026}
X.~Ma and Y.~Zuo,
\emph{The quotient of Milnor number by Tjurina number of hypersurface
singularities in arbitrary characteristic},
arXiv:2604.17757v1 (2026), non-refereed preprint.
\url{https://arxiv.org/abs/2604.17757}.

\bibitem{Parusinski1991}
A.~Parusi\'nski,
\emph{Multiplicity of the dual variety},
Bull. Lond. Math. Soc. \textbf{23} (1991), 429--436.
\href{https://doi.org/10.1112/blms/23.5.429}
{doi:10.1112/blms/23.5.429}.

\bibitem{PhuongTran2023}
H.~V.~N.~Phuong and Q.~H.~Tran,
\emph{A new proof of Stanley's theorem on the strong Lefschetz property},
Colloq. Math. \textbf{173} (2023), 1--8.
\href{https://doi.org/10.4064/cm8987-11-2022}
{doi:10.4064/cm8987-11-2022}.

\bibitem{Stanley1980}
R.~P.~Stanley,
\emph{Weyl groups, the hard Lefschetz theorem, and the Sperner property},
SIAM J. Algebraic Discrete Methods \textbf{1} (1980), 168--184.
\href{https://doi.org/10.1137/0601021}{doi:10.1137/0601021}.

\bibitem{Wahl1985}
J.~M.~Wahl,
\emph{A characterization of quasi-homogeneous Gorenstein surface
singularities},
Compos. Math. \textbf{55} (1985), 269--288, especially Example~4.7.
\url{https://www.numdam.org/item/CM_1985__55_3_269_0/}.

\end{thebibliography}

\end{document}
